\documentclass[11pt,a4paper]{article}
\usepackage{curso}
\begin{document}
\portada{Lección}{10.4}{Transformers dispersos}

\noindent\textbf{Unidad 10.}\quad\textbf{Objetivos:} capas dispersas de grado fijo; comparación a igual presupuesto de conexiones (iso-FLOP); mezcla de información en
grafos de grado fijo.
\textbf{Prerrequisitos:} 5.7, 10.2.

\section{Capas dispersas}
Cada unidad de salida se conecta con $k$ entradas según un patrón fijo (aleatorio, local, por bloques). Costo: $C=n\,k$ multiplicaciones-sumas frente a $n\,d$ de la capa
densa. En un TSU, $k$ lo fija el chip.

\section{Comparación iso-FLOP \evS}
Experimento: $d=64$ entradas, presupuesto $C=8192$, características $\tanh$ dispersas aleatorias, lectura por regresión ridge sobre un maestro no lineal ($N=4000$ de
entrenamiento, 2000 de prueba, 3 máscaras).
\begin{center}\small\begin{tabular}{rrr}\toprule
grado $k$ & ancho $n$ & ECM normalizado\\\midrule
64 (densa) & 128 & 0{,}115\\ 32 & 256 & 0{,}092\\ 16 & 512 & 0{,}083\\ 8 & 1024 & 0{,}091\\ 4 & 2048 & 0{,}101\\ 2 & 4096 & 0{,}109\\\bottomrule
\end{tabular}\end{center}
Curva en U: con grado muy bajo cada unidad ve demasiado poco. El óptimo de este ejemplo no se generaliza sin más a modelos grandes; las leyes de escala dispersas
(Frantar et al.) encuentran que, a igual cómputo, la dispersión ayuda con entrenamiento largo y que el óptimo depende del tamaño y los datos.

\section{Mezcla en grafos de grado fijo}
Campo receptivo tras $\ell$ capas: $r\approx\ell k$ con patrón local, $r\sim k^\ell$ con patrón de mundo pequeño (local más saltos largos). Con $k=16$ y 4096 entradas:
unas 256 capas locales frente a 3 de mundo pequeño.

\begin{aplicacion}
Elegir $k$ según el hardware y ajustar el ancho a presupuesto fijo; máscaras estáticas desde el inicio. Laboratorio L30: barrido de grado en un modelo pequeño.
\end{aplicacion}

\begin{error}
\begin{itemize}[nosep]
\item Comparar con distinto presupuesto. \quad Suponer que lo disperso es rápido en GPU.
\item Patrones solo locales en redes poco profundas. \quad Generalizar el óptimo de un ejemplo pequeño.
\end{itemize}
\end{error}

\begin{thebibliography}{9}\small
\bibitem{frantar} E. Frantar et al., Scaling laws for sparsely-connected foundation models, \emph{ICLR} (2024).
\bibitem{child} R. Child et al., Generating long sequences with sparse transformers, arXiv:1904.10509 (2019).
\bibitem{z1t} Extropic (G. Verdon et al.), Z1T: Sparse transformer-like models for probabilistic hardware, 2026.
\end{thebibliography}
\end{document}
